% !Mode:: "TeX:UTF-8"
\chapter{基于全局预测的专家提前调度方法}
\label{chap:global_schedule}

\section{问题描述与设计动机}
在 VL-MoE 推理系统中，MoE 层的专家访问由 Router 动态决定。对于显存受限的部署场景，专家权重往往无法全部常驻 GPU，而需要在 CPU 与 GPU 之间按需搬运。若系统等到真实访问发生时才开始加载专家，就会在专家计算路径上引入明显等待时间，从而放大 TTFT 和解码阶段时延。

围绕这一问题，本文将其定义为面向 offload 场景的\textbf{专家提前调度问题}：利用当前层 Router 已经产生的路由信号，预测未来若干层更可能激活的专家集合，并在真实访问发生前完成预取、驻留和替换决策，从而把一部分等待时间前移到当前层的计算窗口中。现有系统工作已经从 KV cache、prefill/decoding 解耦以及 expert offloading 等角度探索过 serving 优化\cite{kwon2023pagedattention,zhong2024distserve,xue2024moeinfinity,he2024expertflow,zhu2026dali}。本文进一步利用\textbf{路由可预测性}，把优化对象从“被动处理专家访问”推进到“主动规划专家访问”。

与纯文本 MoE 相比，Qwen3-VL\cite{bai2025qwen3vl} 这类多模态模型还具有额外的条件变量。图像 Token 与文本 Token 的路由密度、位置分布和上下文稳定性并不相同，因此本文在预测器输入和调度器决策中都显式引入模态信息，而不是把它当成后处理噪声。
面向 offload 场景的全局预测与专家调度整体框架如图~\ref{fig:fig4-offload-arch} 所示。

从系统基础设施角度看，MoE offload 的本质是在有限 GPU 显存和较慢 CPU-GPU 数据通道之间做时间换空间。大模型算力基础设施研究指出，随着模型规模和模态复杂度提升，计算、显存、互联和存储之间的协同成为推理部署的重要瓶颈\cite{zhang2024computeinfra}。对于 VL-MoE 而言，这一矛盾更加明显：模型总参数量由大量专家构成，但单个 Token 实际只访问其中一小部分专家。如果能够提前知道未来访问路径，就可以让部分专家权重在计算发生前完成加载，从而缓解同步等待。


\begin{figure}[!htbp]
\centering
\includegraphics[width=0.95\textwidth]{fig4_1_offload_architecture}
\caption{面向 offload 场景的全局预测与专家调度整体框架}
\label{fig:fig4-offload-arch}
\end{figure}

从形式化角度看，设第 $l$ 层输入的 Token 集合为 $\mathcal{X}^{(l)}=\{x_i^{(l)}\}_{i=1}^{n_l}$，Router 为每个 Token 选择 Top-$k$ 专家集合
\begin{equation}
\mathcal{A}_i^{(l)} = \mathrm{TopK}\big(\mathrm{Router}(x_i^{(l)})\big), \qquad \mathcal{A}_i^{(l)} \subseteq \mathcal{E},
\end{equation}
其中 $\mathcal{E}=\{e_1,e_2,\ldots,e_M\}$ 表示全部专家集合，$M$ 为专家总数。
本章的目标是在保证预测额外开销可控的前提下，利用当前层可观测信号逼近未来层的真实激活集合，从而尽可能减少专家加载时延、控制额外数据搬移并降低等待开销。

更具体地说，运行时调度需要回答三个问题：第一，未来哪些层的哪些专家更可能被访问；第二，在 GPU expert cache 容量有限的情况下，哪些候选专家值得提前搬入；第三，当预测失败或缓存压力上升时，如何回退到按需加载而不破坏计算正确性。后续小节分别从数据构建、预测器结构和受限预取策略三个层面回答这些问题。

\section{数据采集与样本构建}
为了让预测器直接对真实 VL 路由模式建模，本文没有复用旧的纯文本 MoE 语料，而是重新采集了 Qwen3-VL 的 fresh 路由 trace。采集链路由三步组成：
\begin{enumerate}
    \item 在 Qwen3-VL 推理过程中对 MoE Router 增加 hook，记录每层的 router 概率、Top-$k$ 专家、Token 位置、输入标识和模态标识；
    \item 将单次请求的逐层路由组织为 token 级样本，保留 layer-0 路由信号以及后续层监督标签；
    \item 使用构建脚本将原始 trace 转换为训练所需的 Parquet 数据集，再切分为训练集和验证集。
\end{enumerate}

在当前实验中，fresh 数据集共包含 8347 个 token，其中训练集 7268 个 token，验证集 1079 个 token；按模态统计，图像 token 为 6734，文本 token 为 1613。数据集覆盖了多项选择、描述、推理和图文组合等样本，因此能够提供较丰富的路由模式。
fresh 路由数据的采集与转换流程如图~\ref{fig:fig4-data-pipeline} 所示，训练数据规模和模态分布统计如图~\ref{fig:fig4-dataset} 所示。

在数据构建过程中，本文特别保留了请求级边界，而不是只把所有 token 打散后随机划分。这是因为同一请求内部的视觉 Token 和文本 Token 具有相同图像、prompt 和上下文条件，如果训练集和验证集同时包含同一请求的 token，验证指标会高估模型对新请求的泛化能力。按请求划分能够更真实地评估预测器是否学到了跨请求可迁移的路由规律。


\begin{figure}[!htbp]
\centering
\includegraphics[width=0.95\textwidth]{fig4_2_data_pipeline}
\caption{fresh 路由数据采集、转换与训练流程}
\label{fig:fig4-data-pipeline}
\end{figure}

\begin{figure}[!htbp]
\centering
\includegraphics[width=0.95\textwidth]{fig4_3_dataset_stats}
\caption{训练数据规模与模态分布统计}
\label{fig:fig4-dataset}
\end{figure}

\section{全局路由预测器设计}
\subsection{输入特征}
本文实现的预测器是一个 layer-0 条件化的全局路由预测器。运行时只需要观察第 0 层 Router 输出，就可以一次性预测后续所有 MoE 层的专家使用情况。最终版本采用带附加特征的配置，即在基础路由信号之外，还引入 Token 位置、输入标识、模态信息和路由统计量。

对于单个 Token，输入特征包括：
\begin{enumerate}
    \item layer-0 router 概率向量 $\mathbf{g}_i^{(0)} \in \mathbb{R}^{M}$；
    \item layer-0 Top-$k$ 专家的一热或 multi-hot 表示 $\mathbf{m}_i^{(0)} \in \{0,1\}^{M}$；
    \item token 位置编号及其 bucket 化嵌入；
    \item input id 的哈希嵌入；
    \item 是否为视觉 token 的模态嵌入；
    \item router 统计量，包括 Top-$k$ 概率质量、最大概率和归一化熵。
\end{enumerate}

将这些特征拼接后，得到单个 token 的全局预测输入
\begin{equation}
\mathbf{f}_i = \big[\mathbf{g}_i^{(0)} \,\|\, \mathbf{m}_i^{(0)} \,\|\, \phi(\text{pos}_i) \,\|\, \psi(\text{input\_id}_i) \,\|\, \eta(\text{modality}_i) \,\|\, \rho(\mathbf{g}_i^{(0)}, \mathbf{m}_i^{(0)})\big],
\end{equation}
其中 $\rho(\cdot)$ 表示由 router 概率与 Top-$k$ one-hot 计算得到的统计特征映射。

上述特征的设计兼顾了三类信息。第一类是路由信息，包括首层 router 概率和 Top-$k$ one-hot，它们直接反映当前 token 在专家空间中的初始偏好；第二类是 token 上下文信息，包括位置、输入 id 和模态标识，用于区分图像 token、文本 token 以及不同请求中的局部结构；第三类是统计信息，包括概率质量、最大概率和熵，用于刻画 router 决策是否集中。当首层分布熵较低时，后续路径往往更稳定；当分布较平坦时，预测结果需要更谨慎地进入调度器。


\subsection{网络结构}
预测器采用“共享骨干 + 多层输出头”的轻量结构。首先将路由信号和 token 信号投影到同一隐空间，再经过多层 MLP backbone 提取稳定特征，最后为每一层设置独立输出头，直接输出该层的专家 logits。对应到实现，核心模块位于原型系统的 \texttt{qwen3vl\_global.py}。其中 \texttt{route\_proj} 负责处理 router 概率和 Top-$k$ one-hot，\texttt{token\_proj} 负责处理 token 位置，\texttt{use\_extra\_features} 打开后再接入 input id、模态和 route statistics。训练过程中 recall 与 loss 的变化如图~\ref{fig:fig4-train-curve} 所示。

\begin{figure}[!htbp]
\centering
\includegraphics[width=0.95\textwidth]{fig4_4_training_curve}
\caption{全局预测器训练过程中的 recall 与 loss 变化}
\label{fig:fig4-train-curve}
\end{figure}

\subsection{训练目标}
由于一个 token 在同一层可能激活多个专家，因此本文将该任务建模为多标签学习问题。设真实标签为 $\mathbf{y}\in\{0,1\}^{H\times M}$，预测 logits 为 $\hat{\mathbf{z}}\in\mathbb{R}^{H\times M}$，则训练目标为
\begin{equation}
\mathcal{L}_{\mathrm{pred}} = \sum_{h=1}^{H} \omega_h \cdot \mathrm{BCEWithLogits}\big(\hat{\mathbf{z}}^{(h)}, \mathbf{y}^{(h)}\big),
\end{equation}
其中 $\omega_h$ 为不同预测视野的权重。训练完成后，模型在验证集上的 best mean recall 达到 0.7219，early/mid/late 三个视野的 recall 也都稳定在较高水平，具体结果见表~\ref{tab:global-recall}。这说明 layer-0 条件化预测已经具备进入运行时调度链路的基础。

本文选择 recall 作为核心离线指标，是因为专家预取更关注“是否覆盖未来真实访问专家”。在 offload 场景下，漏掉一个即将访问的专家会导致同步加载等待，而多预测一个专家则主要表现为额外预取和缓存压力。二者都需要控制，但在候选生成阶段，较高 recall 是形成调度收益的前提。后续调度器再通过置信度过滤和预算限制减少误预测带来的成本。


\begin{table}[!htbp]
\centering
\caption{全局预测器在验证集上的分视野 recall}
\label{tab:global-recall}

\begin{tabular}{lccc}
\hline
指标 & early & mid & late \\
\hline
Recall & 0.7158 & 0.7453 & 0.7042 \\
\hline
\end{tabular}

\end{table}

\section{预测驱动的专家预取与缓存调度}
\subsection{预测驱动的受限预取调度策略}
全局预测器会给出后续多层可能访问的专家集合，但这些预测结果不能被直接等价为搬运指令。一方面，预测本身存在误差，低置信度候选专家如果被提前搬入 GPU，可能只会增加无效传输；另一方面，GPU expert cache 容量有限，过多预取还可能挤出当前层真实计算需要的专家。因此，本文采用预测驱动的受限预取调度策略，将预测结果先压缩为少量高价值候选，再在不破坏当前工作集的前提下执行预取。

具体而言，调度器首先根据全局预测结果形成未来若干层的候选专家集合，并从置信度、层距和预测区分度三个角度为候选专家打分。对于未来第 $h$ 层的候选专家 $e$，调度分数定义为
\begin{equation}
s(e,h) = \alpha_h \cdot p(e,h) \cdot \max(p_1-p_2,0),
\end{equation}
其中 $p(e,h)$ 表示专家 $e$ 在未来第 $h$ 层被访问的预测概率，$p_1$ 与 $p_2$ 分别表示该预测分布中的 top-1 和 top-2 概率，$\alpha_h$ 为层距衰减项。该分数同时体现三层含义：预测概率越高，候选专家越值得预取；距离当前层越近，预取命中的时间价值越高；top-1 与 top-2 的 margin 越大，说明预测越明确，误预取风险越低。

在得到候选分数后，系统继续执行三类约束。第一，通过置信度阈值和最小 margin 过滤掉不稳定预测，避免低价值专家进入预取队列。第二，通过每层 key 上限和全局预取预算限制搬运规模，使额外传输量保持可控。第三，将当前层正在访问的真实专家作为保护集合，在缓存替换时禁止预取候选挤出当前计算必需的专家。

最终，\path{select_weighted_global_prefetch_plan()} 会按照综合分数生成有限长度的预取列表。若候选专家已在 GPU cache 中，系统只更新其缓存分数；否则在预算允许且不破坏保护集合的条件下提前搬入 GPU，将部分同步等待转化为可重叠的后台搬运。

因此，该策略的核心可以概括为“预测候选、筛选高值、保护当前、限量预取”，即在有限 GPU expert cache 中优先保留近未来最可能带来等待时间下降的专家。第五章实验将进一步验证其在真实 CPU-GPU expert offload 场景中的收益。

从调度生命周期看，一次预取决策包含以下步骤。首先，系统收集当前层真实激活专家，并将其加入保护集合，确保当前计算所需专家不会被预取候选替换。其次，预测器根据 layer-0 条件化特征输出未来层候选专家分数，调度器按层距和置信度进行衰减与过滤。再次，调度器检查候选专家是否已经在 GPU cache 中：若已存在，则提高其驻留优先级；若不存在，则在预算允许时发起异步搬运。最后，当真实访问到达时，系统根据命中情况更新缓存分数，并对未命中专家执行同步加载。


\subsection{预测错误与回退机制}
预测驱动的调度不可避免会出现误预测，因此系统需要在正确性与收益之间做折中。本文采用两级回退机制：
\begin{enumerate}
    \item 对未命中的真实专家立即同步加载，保证计算可继续；
    \item 将命中情况反馈到缓存分数，降低低价值专家的长期驻留优先级。
\end{enumerate}

因此，本文的目标不是让预测永远正确，而是让正确预测尽量提前生效，同时把错误预测控制在可接受范围内。

需要强调的是，回退机制保证的是正确性，而不是最优性能。在最坏情况下，如果预测完全失效，系统仍会退化为 demand 策略，即真实访问发生时同步加载专家。这样可以避免预测模块成为系统稳定性的单点风险，也便于在不同部署环境中逐步开启预测式调度。

\section{系统实现与复杂度分析}
整个原型实现拆分为三个阶段：
\begin{enumerate}
    \item 数据采集与构建：\path{collect_qwen3vl_global_scratch_data.py}、\path{build_qwen3vl_global_scratch_parquet.py}；
    \item 预测器训练：\path{train_qwen3vl_global_predictor.py}；
    \item 在线集成：\path{qwen3vl_global.py}、\path{global_scheduling.py}、\path{run_qwen3vl_real_offload_benchmark.py} 与 \path{run_qwen3vl_end_to_end_offload.py}。
\end{enumerate}

从复杂度上看，预测器的主要代价来自 MLP 前向和多层输出头，但其输入维度和参数规模都远小于 MoE 专家权重搬运成本，因此适合作为调度前置模块。调度器本身的复杂度主要来自候选专家排序和缓存替换，均低于专家前向与权重传输。若将 token 数、预测视野和候选专家数分别记为 $B$、$H$ 与 $K$，则预测与筛选的开销可近似写为
\begin{equation}
\mathcal{O}(B \cdot H \cdot K),
\end{equation}
而缓存更新和排序通常可压缩到候选集合规模上的 $O(K\log K)$。

与第三章的视觉侧优化不同，第四章方法的收益主要体现在等待时间重叠而不是算子计算量减少。预测器本身不会减少真实需要执行的专家 MLP 次数，也不一定减少总传输量；它的作用是让高概率专家更早进入 GPU cache，使一部分 CPU-GPU 搬运与当前层计算重叠。因此，评估该方法时不能只看传输 GB 数，还需要同时观察同步加载次数、缓存命中次数和总耗时。

\section{本章小结}
本章基于 fresh 路由数据构建了 layer-0 条件化全局预测器，并设计了预测驱动的受限预取调度策略。该方法将预测信号引入专家预取、缓存驻留和回退机制，使专家搬运从“访问后处理”转向“访问前准备”；下一章将在真实 CPU-GPU expert offload 微基准中验证其收益。
